\chapter{PHÂN TÍCH CÁC CHÍNH SÁCH KHAI THÁC NGHỀ CÁ}
\label{ch:policies}

\textit{(Nội dung chi tiết của chương này sẽ được phát triển đầy đủ trong giai đoạn tiếp theo):}

\section{Chính sách Sản lượng bền vững tối đa (MSY)}
\begin{itemize}
    \item Tối đa hóa hàm sản lượng đánh bắt $Y(E) = q_1 E x^* + q_2 E y^*$.
    \item \textbf{Trường hợp 1 (Chỉ khai thác cá vược sọc):} Nỗ lực $E_{\mathrm{MSY}} = \frac{rd - ae}{2aq_2}$. Chứng minh quần thể kẻ săn mồi ở mức MSY bằng 50\% quần thể khi chưa đánh bắt (Định lý 4.1).
    \item \textbf{Trường hợp 2 (Chỉ khai thác cá mòi dầu):} Điều kiện tồn tại MSY đòi hỏi $\frac{rd - ae}{cr + be} > \frac{d}{2c}$. Đối với hệ sinh thái vịnh Chesapeake, điều kiện này không thỏa mãn, nỗ lực $E_{Y_m} \approx 1.02$ sẽ dẫn tới sự tuyệt chủng của cá vược sọc; do đó chính sách khuyến nghị giảm 50\%: $E_{\mathrm{MSY}} = 0.5$ (Định lý 4.2).
    \item \textbf{Trường hợp 3 (Khai thác đồng thời cả hai loài):} Điều kiện tồn tại nghiệm bền vững và mức khuyến nghị an toàn sinh thái $E_{\mathrm{MSY}} = 0.5$ (Định lý 4.3).
\end{itemize}

\section{Chính sách Sản lượng kinh tế tối đa (MEY)}
\begin{itemize}
    \item Tối đa hóa hàm tô kinh tế ròng $\mathrm{Rent}(E) = R(E) - C(E) = \hat{p} Y(E) - \hat{c} E$.
    \item Chứng minh tính chất quan trọng: $E_{\mathrm{MEY}} < E_{\mathrm{MSY}}$ khi chỉ đánh bắt cá vược sọc, khẳng định chính sách MEY thân thiện hơn với bảo tồn sinh thái.
    \item Phân tích nghịch lý kinh tế khi đánh bắt cá mòi dầu hoặc cả hai loài trong vịnh Chesapeake: $E_{\mathrm{MEY}}$ nhận giá trị âm dưới các giả định chi phí thông thường, chứng minh MEY không thể áp dụng độc lập mà cần kết hợp bảo trợ hệ sinh thái.
\end{itemize}

\section{Chính sách Sản lượng tối đa hóa độ đàn hồi (RMY)}
\begin{itemize}
    \item Đánh giá tác động của nỗ lực khai thác lên độ hồi phục $R$ và thời gian cân bằng $\tau$.
    \item \textbf{Đánh bắt cá vược sọc:} Độ đàn hồi tăng tuyến tính theo nỗ lực trong miền khảo sát ($E_{\mathrm{RMY}} = 10$).
    \item \textbf{Đánh bắt cá mòi dầu:} Độ đàn hồi giảm đơn điệu ngặt theo $E_1$, đạt cực đại tại $E_1 = 0$; bất kỳ nỗ lực đánh bắt cá mòi dầu nào cũng phải trả giá bằng sự suy giảm độ bền vững sinh thái.
    \item \textbf{Đánh bắt kết hợp hai loài:} Sự đánh đổi gay gắt giữa sản lượng và độ ổn định hệ thống.
\end{itemize}
